% chapters/08_fazit.tex

\chapter{Fazit und Ausblick}
\label{chap:fazit}

\section{Fazit}
\label{sec:fazit}

Diese Projektarbeit ist der Frage nachgegangen, ob und wie sich das
Vision-Language-Action-Modell NVIDIA GR00T~N1.6 auf dem Unitree~G1 mit
DEX3-Hand für eine Block-Stacking-Aufgabe feinabstimmen lässt -- und welche
Aussagekraft die anschließende Simulationsevaluation besitzt. Die zentralen
Erkenntnisse:

\begin{itemize}
  \item \textbf{Reproduzierbare Infrastruktur.} Die containerisierte, rein
    über Umgebungsvariablen konfigurierte Pipeline läuft identisch lokal
    (Docker) und auf dem KISSKI-Cluster (Apptainer/\ac{slurm}) und skaliert
    auf vier GPUs.
  \item \textbf{Erfolgreiches Training.} Das selektive Fine-Tuning
    konvergierte über 175\,000~Schritte sauber; der Loss plateaut
    ab~$\approx$~100\,000~Schritten.
  \item \textbf{Domain-Gap als Kernbefund.} Trotz guter
    Open-Loop-Armvorhersagen erreichte die Politik im geschlossenen
    Regelkreis $0/20$~Erfolge und \enquote{fror ein}. Die Dreifach-Diagnose
    (Open-Loop, Replay, Closed-Loop) isolierte die \emph{Beobachtung} als
    einzige Differenz; die SigLIP-Messung quantifizierte den visuellen
    Domain-Gap mit der handnahen Kamera (0,427) als kritischem Blocker.
  \item \textbf{Einordnung statt Fehlinterpretation.} Literaturgestützt
    (SIMPLER~\cite{simpler2024}) ist $0\,\%$ Closed-Loop-Erfolg für ein auf
    Realbildern trainiertes Modell mit eingefrorenem Encoder das
    \emph{erwartete} Ergebnis -- kein Urteil über die Modellqualität. Die
    domain-gap-freie Metrik (Open-Loop-\ac{mse}) spricht dafür, dass das
    Modell die Aufgabe im Prinzip gelernt hat.
  \item \textbf{Extern validierte Pipeline, eingegrenzter Gap.} Die
    Inferenz-Kette reproduzierte einen publizierten Benchmark auf
    0,1~Prozentpunkte (RoboCasa GR-1: 47,7\,\% vs.\ 47,8\,\%,
    \cref{sec:robocasa}); Kamera-Neukalibrierung und Albedo-Korrekturen
    senkten den mittleren Real$\to$Sim-Abstand unter die
    Real$\to$Real-Referenz (\cref{sec:neukalibrierung}) -- dominante
    Gap-Ursache war die Materialfarbe, nicht die Beleuchtung.
  \item \textbf{Verhaltensbeleg statt bloßer Bildähnlichkeit.} Das
    \texttt{span}-Gate (\cref{sec:span-gate}) zeigte: Dieselbe Politik
    kommandiert auf \emph{echten} Bildern die volle Greifbewegung
    (Verhältnis~1,00), in der Simulation nur ein Fünftel. Die Politik hat
    das Greifen gelernt und ruft es nur bei synthetischer Beobachtung nicht
    ab.
  \item \textbf{Checkpoint-Auswahl war messbar falsch.} Der dritte Lauf --
    der erste mit echter Testmenge -- zeigte eine U-Kurve des
    Validierungsfehlers: Der letzte Checkpoint ist um 25\,\% schlechter als
    der beste (\cref{sec:lauf3}); Overfitting ist damit erstmals gemessen
    statt vermutet.
  \item \textbf{RL-Pipeline implementiert, Lernwirkung offen.} Der
    FPO-Trainer absolvierte eine vollständige Trainingsiteration auf
    institutseigener Blackwell-Hardware (\cref{sec:rl-stand}); vor einem
    produktiven Lauf stehen Co-Training (\cref{sec:cotrain}) und die
    wiederholten Basismessungen.
  \item \textbf{Ein widerrufener Fix als methodische Warnung.} Die früher
    als Korrektur verbuchte Vorzeichen-Spiegelung ließ die Hand in allen
    Sim-Läufen unvollständig schließen (\cref{sec:vorzeichen-widerruf});
    die simulationsseitigen Greifzahlen stehen daher unter ausdrücklichem
    Vorbehalt (\cref{sec:vorbehalt-greifzahlen}).
\end{itemize}

Der wichtigste Beitrag der Arbeit ist damit weniger eine hohe Erfolgsrate als
die \emph{methodisch saubere Diagnose}, warum die naheliegende Sim-Evaluation
für dieses Modell die falsche Messgröße ist -- mit der konkreten Empfehlung,
Modellqualität primär über Open-Loop-Eval und reale Roboter-Rollouts zu
belegen (\cref{sec:konsequenz}).

\section{Ausblick}
\label{sec:ausblick}

\paragraph{Messungen mit korrigierter Geometrie wiederholen.}
Der unmittelbarste Schritt folgt aus dem Widerruf des Vorzeichen-Fixes: Die
simulationsseitigen Greifgrößen sind mit korrigierter Hand- und
Basisgeometrie neu zu erheben. Erst danach existiert ein belastbarer
Nullpunkt für Co-Training und RL -- und erst danach lässt sich beziffern,
welcher Anteil des Sim-Versagens tatsächlich auf den visuellen Domain-Gap
entfällt.

\paragraph{Co-Training durchführen.}
Der Lauf auf echten \emph{und} gerenderten Bildern ist der inhaltlich
nächste Schritt: Er erfüllt als einziger die Bedingung aus
\cref{sec:frozen-vs-tuned}, dass der Vision-Encoder im Training
Simulationsbilder sieht -- auf reinen Realdaten verschlechterte das
Encoder-Tuning das Verhalten bis zum Politik-Kollaps
(\cref{sec:vision-lauf}). Der gerenderte Datensatz ist mit der korrigierten
Geometrie neu zu erzeugen; parallel dazu läuft der generative
Cosmos-Transfer-Strang (\cref{sec:cotrain}).

\paragraph{Reinforcement-Learning-Fine-Tuning produktiv machen.}
RL adressiert die \ac{bc}-Schwäche an der Wurzel und löst den Domain-Gap
konzeptionell auf, da in derselben Simulation trainiert und evaluiert wird.
Die FPO-Pipeline ist implementiert und end-to-end validiert
(\cref{sec:rl-stand}); nach Co-Training und Basismessungen stehen der erste
produktive Lauf und der Nachweis der Lernwirkung an.

\paragraph{Lokomotion.}
Mit der entkoppelten Architektur (\cref{sec:loco-pfade}) lässt sich der
bislang fixierte Roboter zum Laufen bringen, ohne neue Manipulationsdaten zu
erheben -- gestuft von der Sim-Lauf-Demo bis, langfristig, zur unified
Whole-Body-\ac{vla}.

\paragraph{Sim-to-Real.}
Der entscheidende ausstehende Test ist der Einsatz auf der physischen
G1-Hardware. Da das Closed-Loop-Versagen ein sim-spezifisches
Wahrnehmungsproblem ist, ist plausibel, dass das Modell auf realen
Kamerabildern -- die der Trainingsverteilung entsprechen -- deutlich besser
abschneidet. Der reale Roboter bleibt der Goldstandard und der nächste große
Meilenstein des Projekts.
